\section{Conclusion}


In this paper, we propose a novel Unified Dual-Mask Physical Model for non-Lambertian light field depth estimation from plenoptic cameras. By decoupling angular signals via physical priors and introducing a geometry-aware adaptive routing mechanism, our framework effectively overcomes the inherent geometric constraint failures of traditional epipolar plane image (EPI) (EPI) analysis under complex reflectance conditions. 

Specifically, the core contributions of this work are summarized as follows:
\textit{ \textbf{Physics-Prior Angular Signal Decomposition and Geometric Material Classification:} We constructed a three-layer angular signal decomposition model that decouples light field signals at both physical and geometric levels. This resolves the resolution insufficiency of traditional spectral methods under $9 \times 9$ discrete angular sampling, providing a reliable geometric foundation for precise non-Lambertian surface identification and establishing the theoretical core of the unified physical model.
} \textbf{Geometric Root Cause Analysis of EPI Failure and Dual-Branch Adaptive Routing:} We explicitly analyzed the physical limitation boundaries of traditional EPI-based methods, thereby avoiding ineffective parameter tuning in loss functions and learning rates. By designing a dual-branch adaptive routing mechanism guided by geometric features, we provided clear theoretical guidance for network architecture design, proving it as an effective paradigm to mitigate depth estimation degradation on non-Lambertian surfaces.
\textbf{Cross-Domain Global Statistical Validation and Multi-Domain Balanced Training:} We validated that the theoretically decomposed features possess strong generalization capabilities across diverse datasets and domains. The proposed multi-domain balanced training framework significantly enhances model robustness under uneven data distributions, providing reliable data and training support for the engineering deployment of unified light field depth estimation.

Finally, since we showed that the current physical decomposition inherently assumes locally smooth surfaces and may struggle with highly intricate micro-structures or severe inter-reflections, future work includes extending our dual-mask routing mechanism to incorporate neural implicit representations. This extension will enable more precise volumetric depth rendering for translucent and participating media in complex real-world light field scenarios.


The proposed Unified Dual-Mask Physical Model introduces a paradigm shift from purely data-driven feature extraction to physics-informed signal decoupling. By explicitly formulating the angular signal as $I(u,v) = DC + a\cdot u + b\cdot v + \epsilon(u,v)$, we address the fundamental aliasing problem inherent in low-resolution plenoptic sampling (e.g., $9 \times 9$). From a signal processing perspective, the failure of traditional 2D-DFT spectral analysis under such sparse angular sampling is an inevitable consequence of violating the Nyquist-Shannon sampling theorem for high-frequency specular components. Our transition to spatial-angular geometric features, such as angular gradients and correlation coherence, bypasses this frequency-domain bottleneck, offering a mathematically robust alternative for material characterization without requiring dense angular arrays.

Beyond depth estimation, the three-layer decomposition model holds significant potential for broader light field computational photography tasks. The isolated residual term $\epsilon(u,v)$, which encapsulates complex non-Lambertian reflections, could serve as a physically meaningful prior for light field reflection removal, specularity editing, and novel view synthesis under extreme illumination conditions. Furthermore, the geometric root cause analysis of Epipolar Plane Image (EPI) degradation provides a theoretical blueprint for designing next-generation multi-view stereo algorithms. It suggests that explicit physical routing is fundamentally superior to implicit attention mechanisms when dealing with severe photo-consistency assumption violations.

Despite these theoretical and empirical advantages, several limitations and engineering trade-offs warrant discussion. First, the dual-branch adaptive routing mechanism, while effective in mitigating EPI structural collapse, introduces additional computational overhead. The random forest classifier and the parallel processing branches increase the inference latency and memory footprint compared to single-branch end-to-end networks. Future engineering optimizations could explore lightweight differentiable routing modules or knowledge distillation to compress the dual-branch architecture for real-time embedded deployments in robotic navigation or autonomous systems. 

Second, the current physical model assumes that the residual term $\epsilon(u,v)$ sufficiently captures all non-Lambertian behaviors. However, for highly transparent objects (e.g., thick glass with complex multi-bounce refractions) or anisotropic materials with direction-dependent BRDFs, the linear approximation of the depth term ($a\cdot u + b\cdot v$) may still suffer from localized geometric distortions. In such extreme scenarios, the geometric features extracted from $\epsilon(u,v)$ might conflate high-frequency sensor noise with genuine volumetric scattering or fine-grained anisotropic highlights, potentially degrading the classification accuracy and subsequent depth regression. Incorporating polarization cues or active illumination could further disentangle these highly complex optical phenomena.

Finally, while our framework excels in explicit geometric reconstruction, the rapid advancement of implicit neural representations, such as Neural Radiance Fields (NeRF) and 3D Gaussian Splatting, presents a promising avenue for future integration. Embedding our physics-prior angular decomposition and dual-mask routing logic into the rendering equations of implicit models could bridge the gap between explicit physical interpretability and the photorealistic rendering capabilities of neural fields, ultimately enabling unified 3D reconstruction and relighting for highly complex, non-Lambertian real-world scenes.

Despite the demonstrated efficacy of the proposed unified physical model, several limitations remain that warrant further investigation. First, the three-layer angular signal decomposition model, formulated as $I(u,v) = DC + a\cdot u + b\cdot v + \varepsilon(u,v)$, inherently assumes a locally constant ambient illumination and linear parallax variation. In scenarios with severe occlusions, high-frequency environmental lighting, or complex inter-reflections, the linear approximation degrades, causing the residual term $\varepsilon(u,v)$ to conflate material reflectance with geometric artifacts. This conflation fundamentally limits the upper bound of depth accuracy in highly cluttered scenes. Furthermore, the current dual-branch routing mechanism relies on confidence-based thresholding, which may introduce boundary artifacts at the transition zones between Lambertian and non-Lambertian regions. Second, while the geometry-aware random forest classifier achieves high accuracy on benchmark datasets, its generalization to unseen, highly anisotropic or complex subsurface scattering materials in the wild may suffer from domain shift. Finally, although our geometric feature extraction circumvents the spectral aliasing of 2D-DFT at $9 \times 9$ angular resolution, the discrete differentiation required for angular gradient computation becomes highly sensitive to quantization noise at even sparser angular samplings (e.g., $5 \times 5$).

To address these challenges, future work will explore several primary directions. First, we plan to extend the linear decomposition into a non-linear physical rendering approximation that explicitly models occlusion boundaries and inter-reflections, thereby purifying the residual reflectance signals. We will also investigate soft-routing mechanisms with spatial regularization to ensure seamless depth transitions across material boundaries. Second, to enhance the robustness of material classification in unconstrained environments, we will explore self-supervised contrastive learning frameworks that leverage real-world, unannotated light fields to align geometric feature distributions across diverse domains. Third, to mitigate discrete differentiation noise at ultra-sparse angular resolutions, we intend to develop continuous angular interpolation modules guided by physical priors. Additionally, incorporating temporal consistency constraints for dynamic light field video sequences presents a promising avenue for robust non-Lambertian tracking. Lastly, we aim to integrate our dual-mask physical priors into emerging neural scene representations, such as 3D Gaussian Splatting. By embedding the physics-based angular decomposition into the differentiable rendering pipeline, we can jointly optimize depth geometry and spatially varying BRDFs, paving the way for holistic 3D reconstruction of complex non-Lambertian scenes.